%=======================================================================
%  Progetto di Reti Logiche - AA 2025/2026
%  Relazione del modulo project_reti_logiche
%  Gestione di una lista ordinata di task in memoria esterna
%=======================================================================
\documentclass[a4paper,12pt]{article}

% --- Pacchetti lingua e codifica ---
\usepackage[utf8]{inputenc}
\usepackage[T1]{fontenc}
\usepackage{amsmath}
\usepackage[italian]{babel}

% --- Margini e geometria ---
\usepackage{geometry}
\geometry{
    margin=2.5cm,
    headheight=14pt,
    includehead,
    includefoot
}

% --- Tipografia ---
\usepackage{microtype}

% --- Grafica ---
\usepackage{graphicx}
\usepackage{xcolor}
\usepackage{tikz}
\usetikzlibrary{automata, positioning, arrows.meta, shapes.geometric, fit, calc, decorations.pathreplacing}

% --- Tabelle e figure ---
\usepackage{booktabs}
\usepackage{float}
\usepackage{caption}
\usepackage{subcaption}

% --- Collegamenti ipertestuali ---
\usepackage{hyperref}
\hypersetup{
    colorlinks=true,
    linkcolor=blue,
    urlcolor=blue,
    citecolor=blue
}

% --- Intestazioni e piè di pagina ---
\usepackage{fancyhdr}
\usepackage{lastpage}
\pagestyle{fancy}
\fancyhf{}
\fancyhead[L]{\small Progetto di Reti Logiche --- AA 2025/2026}
\fancyhead[R]{\small Pagina \thepage{}/\pageref{LastPage}}
\fancyfoot[C]{\small \texttt{project\_reti\_logiche} --- Cacioppo, Carbone}
\renewcommand{\headrulewidth}{0.4pt}
\renewcommand{\footrulewidth}{0.4pt}

% --- Stile prima pagina (senza header/footer fancy) ---
\fancypagestyle{plain}{%
    \fancyhf{}
    \renewcommand{\headrulewidth}{0pt}
    \renewcommand{\footrulewidth}{0pt}
}

% --- Paragrafi ---
\usepackage{parskip}

% --- Colori per il diagramma FSM ---
\definecolor{stateInit}{HTML}{E8F5E9}
\definecolor{stateIdle}{HTML}{FFF3E0}
\definecolor{stateOp00}{HTML}{E3F2FD}
\definecolor{stateOp01}{HTML}{FCE4EC}
\definecolor{stateOp10}{HTML}{F3E5F5}
\definecolor{stateOp11}{HTML}{FFFDE7}
\definecolor{stateDone}{HTML}{EFEBE9}
\definecolor{stateRead}{HTML}{E0F2F1}
\definecolor{stateWrite}{HTML}{FBE9E7}

\begin{document}

% --- Pagina del titolo ---
\begin{titlepage}
\centering
\vspace*{1cm}

\includegraphics[width=0.35\textwidth]{assets/logo_polimi.png}\\[1cm]

{\large Scuola di Ingegneria Industriale e dell'Informazione}\\[0.5cm]
{\large Corso di Reti Logiche --- Prova Finale}\\[2cm]

{\huge\bfseries Progetto di Reti Logiche}\\[0.4cm]
{\Large Modulo \texttt{project\_reti\_logiche}}\\[0.4cm]
{\Large Gestione di una lista ordinata di task}\\[0.2cm]
{\Large in memoria esterna}\\[1.5cm]

\vfill

\begin{tabular}{rl}
\textbf{Studente:} & Andrea Cacioppo \\
\textbf{Matricola:} & 10902819 \\[0.3cm]
\textbf{Studente:} & Andrea Carbone \\
\textbf{Matricola:} & 10907106 \\
\end{tabular}
\\[1cm]

{\large Anno Accademico 2025/2026}

\vspace*{1cm}
\end{titlepage}

% --- Indice ---
\tableofcontents
\newpage

%=======================================================================
\section{Introduzione}
%=======================================================================

In questo progetto abbiamo sviluppato un modulo VHDL che si occupa di gestire una lista ordinata di task, con la lista che sta in una memoria esterna. Più precisamente, la memoria è una Single-Port Block RAM configurata in modalità Write-First. Le cose che il modulo deve saper fare sono quattro: alzare la priorità di tutti i task presenti, togliere il primo elemento della lista, inserire un task nuovo rispettando l'ordine di priorità, e cancellare tutto.

Il progetto si basa sulle specifiche del documento \emph{PFRL\_Specifiche\_2526\_V2.pdf} del corso di Reti Logiche (Prof.~Fornaciari, A.A.~2025/2026).

Per capire meglio come funziona, facciamo un esempio. Prendiamo una lista con due task: $[\texttt{Task}_1(P=0),\ \texttt{Task}_2(P=2)]$, $\texttt{count}=2$. Se vogliamo inserire (OP~10) un nuovo $\texttt{Task}_3(P=1)$, il nostro modulo parte dall'indirizzo 1 e scorre la memoria finché non incontra un task con priorità maggiore di 1 (che in questo caso è $\texttt{Task}_2$). A questo punto sposta $\texttt{Task}_2$ di una posizione verso destra (si vede bene in Figura~\ref{fig:mem_struct}) e mette $\texttt{Task}_3$ all'indirizzo 2. Il risultato è $[\texttt{Task}_1(P=0),\ \texttt{Task}_3(P=1),\ \texttt{Task}_2(P=2)]$, e il contatore sale a 3. Se invece usiamo l'incremento priorità (OP~00), ogni task sale di 1 (ma non oltre 3), e otteniamo $[\texttt{Task}_1(P=1),\ \texttt{Task}_2(P=3)]$.

La Figura~\ref{fig:mem_struct} mostra come sono organizzati i dati in memoria, mentre la Figura~\ref{fig:blocchi} dà una vista d'insieme dell'architettura.

\begin{figure}[H]
\centering
\begin{tikzpicture}[
    box/.style={draw, minimum width=7cm, minimum height=0.6cm, font=\footnotesize},
    brace/.style={thick}
]
    \node[box, fill=gray!10] (addr0) {Indirizzo 0: \texttt{count} (6 bit)};
    \node[box, below=0pt of addr0] (addr1) {Indirizzo 1: \texttt{ID\_TASK} (6 bit) $\mid$ \texttt{PRIORITY} (2 bit)};
    \node[box, below=0pt of addr1] (addr2) {Indirizzo 2: \texttt{ID\_TASK} (6 bit) $\mid$ \texttt{PRIORITY} (2 bit)};
    \node[box, below=0pt of addr2] (dots) {$\cdots$};
    \node[box, below=0pt of dots] (addrn) {Indirizzo $N$: \texttt{ID\_TASK} (6 bit) $\mid$ \texttt{PRIORITY} (2 bit)};
    \draw[brace, decorate, decoration={brace, mirror}] ([xshift=-0.3cm]addrn.south west) -- ([xshift=-0.3cm]addr0.north west)
        node[midway, left, font=\footnotesize] {Block RAM};
\end{tikzpicture}
\caption{La memoria: indirizzo 0 = contatore task, indirizzi 1..$N$ = singoli task (6 bit ID, 2 bit priorità).}
\label{fig:mem_struct}
\end{figure}

Abbiamo sintetizzato il tutto con Xilinx Vivado, e passa tutte le simulazioni: pre-sintesi, post-sintesi funzionale e post-sintesi timing (con back-annotation SDF).

\newpage
\subsection{Struttura della memoria}

La memoria è organizzata in modo semplice:

\begin{itemize}
    \item \textbf{Indirizzo 0}: qui si trova il numero di task presenti (count), su 6 bit (quindi da 0 a 63).
    \item \textbf{Indirizzi 1..N}: ogni posizione contiene un task, che occupa un byte. I 6 bit più significativi sono l'\texttt{ID\_TASK}, i 2 bit meno significativi sono la \texttt{PRIORITY}.
\end{itemize}

Per esempio, il byte \texttt{0b00000101} vuol dire ID = 1, priorità = 1. Invece \texttt{0b00001100} è ID = 3, priorità = 0.

La priorità va da 0 (massima) a 3 (minima). I task hanno tutti ID diversi, quindi al massimo ci possono stare 63 elementi. ID = 0 vuol dire lista vuota o errore.

\subsection{Operazioni disponibili}

Con \texttt{i\_op} a 2 bit possiamo scegliere tra quattro operazioni:

\begin{description}
    \item[OP = 00 (Incrementa priorità)] Prende tutti i task in lista e gli alza la priorità di 1, senza superare 3. L'ordine tra i task non cambia: se due finiscono con la stessa priorità, quello che stava prima rimane prima.
    \item[OP = 01 (Rimuovi)] Toglie il primo task (quello all'indirizzo 1). Gli altri si spostano a sinistra per coprire il buco, il contatore scende di 1, e in uscita si vede l'ID di quello rimosso. Se la lista è vuota, il modulo restituisce ID = 0.
    \item[OP = 10 (Inserisci)] Aggiunge un nuovo task, mettendolo in coda a quelli con la stessa priorità (funziona come una FIFO). Se la lista è vuota, finisce direttamente in prima posizione.
    \item[OP = 11 (Clear)] Resetta tutto, contatore a zero.
\end{description}

\newpage
\subsection{Protocollo START-DONE}

Per sincronizzarci con la logica esterna usiamo un semplice handshake START-DONE:

\begin{enumerate}
    \item L'esterno alza \texttt{i\_start} a 1 per far partire l'operazione.
    \item \texttt{i\_start} resta a 1 finché il modulo non risponde alzando \texttt{o\_done} a 1.
    \item Quando vede \texttt{o\_done = 1}, l'esterno può abbassare \texttt{i\_start} a 0.
    \item Il modulo vede lo 0 su \texttt{i\_start} e riporta \texttt{o\_done} a 0, pronto per il prossimo giro.
    \item Una nuova operazione può partire solo quando \texttt{o\_done} è tornato a 0.
\end{enumerate}

Con il reset asincrono (\texttt{i\_rst}), il modulo scrive 0 all'indirizzo 0 della memoria, tiene \texttt{o\_done} a 1 durante il setup, e quando ha finito lo riabbassa.

\subsection{Un esempio pratico}

Vediamo un caso concreto. Partiamo con 4 task in memoria:

\[
\begin{array}{c|c|c}
\text{Indirizzo} & \text{Byte} & \text{Task} \\ \hline
0 & 0\texttt{x}04 & \text{count = 4} \\
1 & 0\texttt{x}40 & \text{ID = 16, P = 0} \\
2 & 0\texttt{x}21 & \text{ID = 8, P = 1}  \\
3 & 0\texttt{x}0E & \text{ID = 3, P = 2}  \\
4 & 0\texttt{x}33 & \text{ID = 12, P = 3}
\end{array}
\]

L'ordine è giusto: P = 0 all'indirizzo 1, poi P = 1, P = 2, e infine P = 3.

\textbf{Inseriamo un nuovo task (OP = 10)}: abbiamo ID = 5, P = 1. Dobbiamo metterlo in coda a quelli con priorità 1, quindi dopo l'indirizzo 2 ma prima dell'indirizzo 3. \newpage Ecco cosa succede:

\[
\begin{array}{c|c|c}
\text{Indirizzo} & \text{Byte} & \text{Task} \\ \hline
0 & 0\texttt{x}05 & \text{count = 5} \\
1 & 0\texttt{x}40 & \text{ID = 16, P = 0} \\
2 & 0\texttt{x}21 & \text{ID = 8, P = 1}  \\
3 & 0\texttt{x}15 & \text{ID = 5, P = 1}  \\
4 & 0\texttt{x}0E & \text{ID = 3, P = 2}  \\
5 & 0\texttt{x}33 & \text{ID = 12, P = 3}
\end{array}
\]

I task con priorità maggiore di 1 sono slittati in avanti di una posizione, e il contatore è passato a 5.

\newpage
%=======================================================================
\section{Architettura del modulo}
%=======================================================================

Il cuore del progetto è una macchina a stati finiti (FSM) con 26 stati, che gestisce tutto il flusso di lavoro. L'abbiamo implementata in un unico processo sincrono (fronte di salita del clock, reset asincrono), perché è il modo più semplice per tenere tutto sotto controllo ed evitare latch indesiderati.

\subsection{Schema funzionale}

\begin{figure}[H]
\centering
\begin{tikzpicture}[
    node distance=1.5cm,
    >=Stealth,
    block/.style={
        draw, rectangle, rounded corners, thick,
        minimum width=3.5cm, minimum height=1.6cm,
        align=center, font=\small
    },
    interface/.style={
        draw, rectangle, rounded corners, thick,
        fill=gray!10,
        minimum width=3.5cm, minimum height=1.4cm,
        align=center, font=\small
    }
]

% Nodi principali
\node[block, fill=stateOp00] (fsm)
    {Control Unit\\FSM (26 stati)};

\node[block, fill=stateRead, right=3cm of fsm] (mem)
    {Memoria\\Block RAM\\Write-First};

\node[block, fill=stateIdle, below=2.5cm of fsm] (regs)
    {Registri e Datapath\\
    count, idx, op, task\_id};

\node[interface, left=3cm of fsm] (input)
    {Ingressi\\
    $i\_clk,\ i\_rst,\ i\_start$\\
    $i\_task\_id,\ i\_op$\\
    $i\_task\_priority$};

\node[interface, above=2cm of mem] (output)
    {Uscite\\
    $o\_done,\ o\_task\_id$};

% Connessioni
\draw[->] (input) -- node[above] {Comandi} (fsm);

\draw[->] (fsm) -- node[above] {\scriptsize addr, data, en, we} (mem);

\draw[->] (mem) -- node[below] {\scriptsize $i\_mem\_data$} (fsm);

\draw[->] (fsm.east) |- node[above, pos=0.63] {$o\_done$} (output.west);

\draw[<->] (fsm) -- node[right] {Controllo e dati} (regs);

% Contorno del sistema
\node[
    draw,
    dashed,
    rounded corners,
    fit={(fsm) (regs)},
    inner sep=0.7cm,
    label={[font=\small]above:project\_reti\_logiche}
] {};

\end{tikzpicture}

\caption{Diagramma a blocchi che illustra l'architettura interna del modulo.}
\label{fig:blocchi}
\end{figure}

\subsection{Modulo 1: Control Unit}

La Control Unit è una FSM a 26 stati che gestisce tutta la logica di controllo: riceve \texttt{i\_start}, campiona gli ingressi, legge il contatore dei task dalla memoria, e decide che sequenza di stati seguire per parlare con la RAM.
\newpage
\subsubsection{Perché un unico processo?}

Mettere tutta la Control Unit in un solo processo VHDL ha tre vantaggi:

\begin{itemize}
    \item \textbf{Meno roba che può rompersi}: transizioni e uscite stanno insieme, niente segnali che girano tra processi diversi.
    \item \textbf{Niente race conditions}: ogni scrittura in memoria capita sul fronte di clock, non ci sono sorprese.
    \item \textbf{Codice più leggibile}: tutto quello che fa la FSM sta in un posto solo, si capisce subito cosa succede.
\end{itemize}

\subsubsection{Gli stati}

I 26 stati della FSM sono organizzati in fasi logiche (Figura~\ref{fig:fsm}). Ecco come si muove il sistema:

\begin{description}
    \item[Inizializzazione --- \texttt{INIT} $\rightarrow$ \texttt{INIT\_WAIT} $\rightarrow$ \texttt{IDLE}.]
          Al reset, \texttt{INIT} scrive 0 all'indirizzo 0. \texttt{INIT\_WAIT} è una specie di sala d'attesa: resta fermo finché \texttt{i\_start} è a 1, e quando scende a 0 passa allo stato di riposo \texttt{IDLE}.

    \item[Idle --- \texttt{IDLE} $\rightarrow$ \texttt{READ\_COUNT}.]
          In \texttt{IDLE} campioniamo gli ingressi (\texttt{i\_op}, \texttt{i\_task\_id}, \texttt{i\_task\_priority}) e aspettiamo \texttt{i\_start}=1. Quando arriva, si va a \texttt{READ\_COUNT}.

    \item[Lettura e decisione --- \texttt{READ\_COUNT} $\rightarrow$ \texttt{WAIT\_COUNT} $\rightarrow$ \texttt{DECODE}.]
          \texttt{READ\_COUNT} punta l'indirizzo 0 con enable alto; \texttt{WAIT\_COUNT} aspetta un ciclo per il dato valido; \texttt{DECODE} prende il conteggio letto e l'operazione e decide la strada.

    \item[Smistamento --- \texttt{DECODE}.] 
          Se il contatore è 0 e l'operazione non ha senso su una lista vuota (incremento o rimozione), salta direttamente a \texttt{DONE}. Altrimenti imbocca il ramo giusto.

    \item[OP 00: Incrementa priorità (3 stati).]
          \texttt{INC\_READ} $\rightarrow$ \texttt{WAIT\_INC} $\rightarrow$ \texttt{INC\_WRITE}. Un loop: per ogni task si legge il byte, si alza la priorità (max 3), si riscrive. Finito, si va a \texttt{DONE}.
          
    \item[OP 01: Rimozione (7 stati).]
          \texttt{RM\_FIRST} prende l'ID del primo task e lo mette in uscita; \texttt{WAIT\_RM\_FIRST} e \texttt{RM\_FIRST\_RD} completano la lettura. Se c'era un solo elemento, si aggiorna il contatore. Altrimenti parte la catena \texttt{SHIFT\_READ} $\rightarrow$ \texttt{WAIT\_SHIFT} $\rightarrow$ \texttt{SHIFT\_WRITE} che sposta ogni task a sinistra per coprire il buco. Poi si decrementa il contatore e si va in \texttt{DONE}.

    \item[OP 10: Inserimento (8 stati).]
          \texttt{FIND} $\rightarrow$ \texttt{WAIT\_FIND} $\rightarrow$ \texttt{DECIDE} scorrono la lista per trovare il primo task con priorità peggiore di quella da inserire. Se serve, si fa lo shift a destra (\texttt{SHIFT\_READ} $\rightarrow$ \texttt{WAIT\_SHIFT} $\rightarrow$ \texttt{SHIFT\_WRITE}) per fare spazio. Poi \texttt{WRITE\_TASK} inserisce, \texttt{UPD\_COUNT} incrementa, fine.

    \item[OP 11: Clear (1 stato).]
          \texttt{CLEAR} scrive 0 all'indirizzo 0 e passa a \texttt{DONE}.

    \item[Fine --- \texttt{DONE}.]
          Alza \texttt{o\_done} a 1. Aspetta che l'esterno abbassi \texttt{i\_start} a 0, poi torna in \texttt{IDLE}.
\end{description}

\begin{figure}[H]
\centering
\resizebox{0.8\textwidth}{!}{% 
\begin{tikzpicture}[
    scale=0.85, transform shape,
    state/.style={
        draw, rounded corners=2pt,
        minimum width=2.2cm, minimum height=0.9cm,
        align=center, font=\footnotesize\bfseries, thick
    },
    edge/.style={
        draw, thick, -{Stealth[scale=0.8]},
        font=\scriptsize, rounded corners=4pt
    },
    group/.style={
        draw, densely dotted, rounded corners=3pt,
        inner sep=6pt,
        font=\footnotesize\itshape
    }
]

% ================= INIT =================
\node[state, fill=stateInit] at (0, 2) (INIT) {INIT};
\node[state, fill=stateInit] at (3.5, 2) (INIT_WAIT) {INIT\_WAIT};

\draw[edge] (INIT) -- (INIT_WAIT);
\path[edge] (INIT_WAIT) edge[loop above] node {start=1} ();
% Percorso pulito e non sovrapposto verso IDLE
\draw[edge] (INIT_WAIT.south) -- (3.5, 1) -- node[above] {start=0} (0, 1) -- (IDLE.north);

% ================= IDLE-READ =================
\node[state, fill=stateIdle] at (0, 0) (IDLE) {IDLE};
\node[state, fill=stateRead] at (3.5, 0) (READ_COUNT) {READ\\COUNT};
\node[state, fill=stateRead] at (7.0, 0) (WAIT_COUNT) {WAIT\\COUNT};
\node[state, fill=stateRead] at (10.5, 0) (DECODE) {DECODE};
\node[state, fill=stateDone] at (16.5, 0) (DONE) {DONE};

\draw[edge] (IDLE.east) -- node[above] {start=1} (READ_COUNT.west);
\draw[edge] (READ_COUNT.east) -- (WAIT_COUNT.west);
\draw[edge] (WAIT_COUNT.east) -- (DECODE.west);

% ================= BUS PRINCIPALE DA DECODE =================
\draw[draw, thick, rounded corners=4pt] (DECODE.south) -- (10.5, -1.5) -- (1.5, -1.5) -- (1.5, -16.5);

% ================= OP 00 =================
\node[state, fill=stateOp00] at (3.5, -3) (INC_READ) {INC\\READ};
\node[state, fill=stateOp00] at (7.0, -3) (WAIT_INC) {WAIT\\INC};
\node[state, fill=stateOp00] at (10.5, -3) (INC_WRITE) {INC\\WRITE};

\draw[edge] (1.5, -3) -- node[above] {op:00} (INC_READ.west);
\draw[edge] (INC_READ.east) -- (WAIT_INC.west);
\draw[edge] (WAIT_INC.east) -- (INC_WRITE.west);
\draw[edge] (INC_WRITE.south) -- (10.5, -4.5) -- node[below] {idx<count} (3.5, -4.5) -- (INC_READ.south);
\draw[edge] (INC_WRITE.east) -- node[above] {idx=count} (16.5, -3) -- (DONE.south);

\node[group, fit={(INC_READ)(WAIT_INC)(INC_WRITE)}, label=above:\textbf{OP 00}] {};

% ================= OP 01 =================
\node[state, fill=stateOp01] at (3.5, -6) (RM1) {RM\\FIRST};
\node[state, fill=stateOp01] at (7.0, -6) (RM2) {WAIT\\FIRST};
\node[state, fill=stateOp01] at (10.5, -6) (RM3) {RM\\FIRST\_RD};
\node[state, fill=stateOp01] at (14.0, -6) (UPD) {UPD\\COUNT};

\node[state, fill=stateOp01] at (3.5, -8.5) (SHIFT_R) {SHIFT\\READ};
\node[state, fill=stateOp01] at (7.0, -8.5) (WAIT_RS) {WAIT\\SHIFT};
\node[state, fill=stateOp01] at (10.5, -8.5) (SHIFT_W) {SHIFT\\WRITE};

\draw[edge] (1.5, -6) -- node[above] {op:01} (RM1.west);
\draw[edge] (RM1.east) -- (RM2.west);
\draw[edge] (RM2.east) -- (RM3.west);
\draw[edge] (RM3.south) -- (10.5, -7.3) -- node[above] {count>1} (2.0, -7.3) |- (SHIFT_R.west);
\draw[edge] (RM3.east) -- node[above] {count=1} (UPD.west);

\draw[edge] (SHIFT_R.east) -- (WAIT_RS.west);
\draw[edge] (WAIT_RS.east) -- (SHIFT_W.west);
\draw[edge] (SHIFT_W.south) -- (10.5, -10.0) -- node[midway, below=4pt] {idx<count} (3.5, -10.0) -- (SHIFT_R.south);
\draw[edge] (SHIFT_W.east) -- node[above] {idx=count} (14.0, -8.5) -- (UPD.south);

\draw[edge] (UPD.east) -- node[above] {fine} (16.5, -6) -- (DONE.south);

\node[group, fit={(RM1)(RM2)(RM3)(SHIFT_R)(WAIT_RS)(SHIFT_W)(UPD)}, label=above:\textbf{OP 01}] {};

% ================= OP 10 =================
\node[state, fill=stateOp10] at (3.5, -11.5) (FIND) {FIND};
\node[state, fill=stateOp10] at (7.0, -11.5) (WAITF) {WAIT\\FIND};
\node[state, fill=stateOp10] at (10.5, -11.5) (DEC) {DECIDE};
\node[state, fill=stateOp10] at (14.0, -11.5) (WRITE_T) {WRITE\\TASK};

\node[state, fill=stateOp10] at (3.5, -14) (SHIFT_I) {SHIFT\\READ};
\node[state, fill=stateOp10] at (7.0, -14) (WAIT_IS) {WAIT\\SHIFT};
\node[state, fill=stateOp10] at (10.5, -14) (SHIFT_IW) {SHIFT\\WRITE};
\node[state, fill=stateOp10] at (14.0, -14) (UPD_I) {UPD\\COUNT};

\draw[edge] (1.5, -11.5) -- node[above] {op:10} (FIND.west);
\draw[edge] (FIND.east) -- (WAITF.west);
\draw[edge] (WAITF.east) -- (DEC.west);
\draw[edge] (DEC.north) -- (10.5, -10.2) -- node[midway, above=4pt] {prio<reg} (3.5, -10.2) -- (FIND.north);
\draw[edge] (DEC.east) -- node[above] {fine} (WRITE_T.west);
\draw[edge] (DEC.south) -- (10.5, -12.5) -- node[above] {shift} (2.0, -12.5) |- (SHIFT_I.west);

\draw[edge] (SHIFT_I.east) -- (WAIT_IS.west);
\draw[edge] (WAIT_IS.east) -- (SHIFT_IW.west);
\draw[edge] (SHIFT_IW.south) -- (10.5, -15.5) -- node[below] {idx>pos} (3.5, -15.5) -- (SHIFT_I.south);
\draw[edge] (SHIFT_IW.north) -- (10.5, -12.5) -- node[above] {idx=pos} (13.8, -12.5) -- ([xshift=-0.2cm]WRITE_T.south);
\draw[edge] ([xshift=0.2cm]WRITE_T.south) -- ([xshift=0.2cm]UPD_I.north);
\draw[edge] (UPD_I.east) -- node[above] {fine} (16.5, -14) -- (DONE.south);

\node[group, fit={(FIND)(WAITF)(DEC)(SHIFT_I)(WAIT_IS)(SHIFT_IW)(WRITE_T)(UPD_I)}, label=above:\textbf{OP 10}] {};

% ================= OP 11 =================
\node[state, fill=stateOp11] at (3.5, -16.5) (CLEAR) {CLEAR};
\draw[edge] (1.5, -16.5) -- node[above] {op:11} (CLEAR.west);
\draw[edge] (CLEAR.east) -- node[above] {fine} (16.5, -16.5) -- (DONE.south);

% ================= BACK TO IDLE =================
\draw[edge] (DONE.north) -- (16.5, 3.5) -- node[above] {start=0} (-1.5, 3.5) -- (-1.5, 0) -- (IDLE.west);

\end{tikzpicture}%
}
\caption{Diagramma degli stati della Control Unit.}
\label{fig:fsm}
\end{figure}

\subsection{Modulo 2: Registri e datapath}

I registri tengono lo stato interno del sistema: contatore dei task, indice di scorrimento, operazione in corso, ID e priorità. La logica di contorno si occupa di saturare le priorità, confrontare valori e calcolare gli shift.

La Tabella~\ref{tab:registri} riassume i registri usati.

\begin{table}[H]
\centering
\begin{tabular}{llp{7cm}}
\toprule
Registro & Dimensione & Scopo \\
\midrule
\texttt{reg\_count}       & 6 bit & Numero di task attualmente in lista \\
\texttt{reg\_idx}         & 6 bit & Indice utile per scorrere o sfalsare la memoria \\
\texttt{reg\_out\_id}     & 6 bit & ID del task appena rimosso (usato in OP 01) \\
\texttt{reg\_insert\_pos} & 6 bit & Posizione calcolata per un nuovo inserimento (OP 10) \\
\texttt{reg\_op}          & 2 bit & L'operazione richiesta (catturata in \newline IDLE) \\
\texttt{reg\_task\_id}    & 6 bit & L'ID del task da inserire (OP 10) \\
\texttt{reg\_task\_priority} & 2 bit & La priorità del task da inserire (OP 10) \\
\bottomrule
\end{tabular}
\caption{Registri interni.}
\label{tab:registri}
\end{table}

\subsection{Modulo 3: Memoria (Block RAM)}

La memoria è una Single-Port Block RAM in modalità Write-First, fornita già pronta dai docenti. Noi abbiamo solo dovuto interfacciarci, rispettando il suo comportamento (descritto nella documentazione Xilinx Vivado):

\begin{itemize}
    \item \textbf{Scrittura}: con \texttt{o\_mem\_en} e \texttt{o\_mem\_we} alti, al clock il dato su \texttt{o\_mem\_data} viene scritto all'indirizzo \texttt{o\_mem\_addr}. Con la modalità Write-First, il dato è subito disponibile anche in uscita su \texttt{i\_mem\_data} dopo 2~ns.
    \item \textbf{Lettura}: \texttt{o\_mem\_en}=1, \texttt{o\_mem\_we}=0. Si legge l'indirizzo e il dato arriva su \texttt{i\_mem\_data} dopo 2~ns.
    \item \textbf{Riposo}: \texttt{o\_mem\_en}=0, la memoria non fa niente e \texttt{i\_mem\_data} tiene l'ultimo valore.
\end{itemize}

Il nostro modulo lavora in due cicli: nel primo imposta indirizzo e enable, nel secondo prende il dato valido.

\subsection{Modulo 4: Ingressi}

Il modulo prende in ingresso:

\begin{itemize}
    \item \texttt{i\_clk}: clock (periodo 20~ns).
    \item \texttt{i\_rst}: reset asincrono, attivo alto, azzera registri e manda in \texttt{INIT}.
    \item \texttt{i\_start}: avvia l'operazione (protocollo START-DONE).
    \item \texttt{i\_op} (2 bit): seleziona l'operazione.
    \item \texttt{i\_task\_id} (6 bit): ID del nuovo task (per OP~10).
    \item \texttt{i\_task\_priority} (2 bit): priorità del nuovo task (per OP~10).
\end{itemize}

Per evitare che i segnali cambino mentre la FSM lavora, copiamo \texttt{i\_op}, \texttt{i\_task\_id} e \texttt{i\_task\_priority} nei registri interni appena entrati in \texttt{IDLE}.

\subsection{Modulo 5: Uscite}

Le uscite del modulo:

\begin{itemize}
    \item \texttt{o\_done}: vale 1 quando l'operazione è finita, torna a 0 quando l'esterno molla \texttt{i\_start}.
    \item \texttt{o\_task\_id} (6 bit): ID del task rimosso (OP~01). Se lista vuota, vale 0.
    \item \texttt{o\_mem\_addr} (16 bit): indirizzo RAM puntato.
    \item \texttt{o\_mem\_data} (8 bit): dato da scrivere (quando \texttt{o\_mem\_we = 1}).
    \item \texttt{o\_mem\_en}: enable della memoria.
    \item \texttt{o\_mem\_we}: write enable.
\end{itemize}

All'inizio di ogni ciclo riportiamo i segnali di controllo memoria ai valori di default, per evitare latch o scritture accidentali.

\subsection{Scelte progettuali}

\subsubsection{Segnali di memoria}

All'inizio di ogni ciclo, prima del \emph{case}, riportiamo \texttt{o\_mem\_en}, \texttt{o\_mem\_we}, \texttt{o\_mem\_addr} e \texttt{o\_mem\_data} ai valori predefiniti. In questo modo evitiamo latch e scritture non volute.

\subsubsection{Shift dei task}

Per la rimozione (OP 01) partiamo dal secondo task e spostiamo tutto a sinistra fino alla fine. Per l'inserimento (OP 10) facciamo il contrario: partiamo dall'ultimo e spostiamo tutto a destra fino al punto giusto. Così nessun dato viene sovrascritto prima di essere stato spostato.

\subsubsection{Casi limite}

\begin{itemize}
    \item \textbf{Lista vuota}: se \texttt{reg\_count} = 0 e chiediamo OP 00 o OP 01, la FSM va direttamente in \texttt{DONE}. Per OP 01, \texttt{o\_task\_id} = 0.
    \item \textbf{Priorità già al minimo}: con OP 00 la priorità si ferma a 3, senza sforare.
    \item \textbf{Stessa priorità in inserimento}: il nuovo task va in coda a quelli con la stessa priorità (FIFO).
    \item \textbf{Lista piena}: massimo 63 task (gli ID sono univoci). Al 64°, la FSM ignora OP~10 e va direttamente in \texttt{DONE}.
\end{itemize}

\newpage
%=======================================================================
\section{Risultati}
%=======================================================================

\subsection{Come abbiamo testato il sistema}

Abbiamo simulato il modulo con \texttt{GHDL} su macOS ARM64 (VHDL-2008, flag \texttt{--std=08 -frelaxed}). Oltre al test base fornito dai docenti, abbiamo scritto un nostro test bench (\texttt{tb\_edge\_cases.vhd}) con 27 scenari divisi in tre gruppi.

Il test bench confronta da solo le uscite della FSM con i risultati attesi. Se qualcosa non quadra, alza una \texttt{failure} e ferma la simulazione.

\subsection{Test per gli Edge Cases}

\subsubsection{Test Base}

Qui abbiamo visto cosa succede con lista vuota o reset improvviso.

\begin{table}[H]
\centering
\small
\begin{tabular}{@{}llp{8cm}@{}}
\toprule
\textbf{Test} & \textbf{Op} & \textbf{Condizione analizzata e risultato ottenuto} \\
\midrule
1  & OP 00 & Lista vuota: il sistema ignora il comando e il conteggio resta 0. \\
2  & OP 01 & Lista vuota: ci restituisce un ID\_TASK pari a 0. \\
3  & OP 10 & Lista vuota: il task viene inserito all'indirizzo 1. \\
4  & OP 11 & Lista vuota: il clear mantiene il contatore a 0. \\
\addlinespace
5  & OP 01 & Unico task presente: lo rimuove e azzera il contatore. \\
6  & OP 00 & Unico task a priorità 0: lo porta correttamente a priorità 1. \\
7  & OP 00 & Unico task a priorità 3: lo satura senza errori. \\
8  & OP 01+10 & Singolo task: viene rimosso e sostituito, alla fine count è 1. \\
\addlinespace
24 & OP 00 + rst & Reset piazzato a metà di un incremento su 10 task: count azzerato. \\
25 & OP 10 + rst & Reset mentre tenta di inserire un task su una lista di 5: count azzerato. \\
\bottomrule
\end{tabular}
\caption{Test base: lista vuota e reset.}
\label{tab:test_base}
\end{table}

\newpage
\subsubsection{Test Funzionali}

Qui abbiamo messo alla prova la logica di controllo: saturazioni multiple, priorità incrociate, shift su liste grandi.

\begin{table}[H]
\centering
\small
\begin{tabular}{@{}llp{8cm}@{}}
\toprule
\textbf{Test} & \textbf{Op} & \textbf{Condizione analizzata e risultato ottenuto} \\
\midrule
9  & OP 00 & Array con 4 task già a priorità 3: nessuno subisce modifiche. \\
10 & OP 00 & Mix casuale P=0,1,2,3 $\rightarrow$ 1,2,3,3: saturano solo i valori giusti e l'ordine regge. \\
11 & OP 00 & Sequenza di tre incrementi su (0,0,1): spinti progressivamente fino a 3. \\
\addlinespace
12 & OP 10 & Inseriamo P=0 in una lista popolata: finisce in testa facendo spostare tutti. \\
13 & OP 10 & Inseriamo P=3 in una lista popolata: viene inserito in coda alla lista. \\
14 & OP 10 & Priorità già esistente: va in coda al suo sottogruppo. \\
15 & OP 10 & 4 insert con stessa priorità: mantiene intatto l'ordine cronologico d'arrivo. \\
\addlinespace
16 & OP 01 & Leviamo il capofila da una lista di 30: 29 elementi scalano a sinistra. \\
17 & OP 10 & Nuovo innesto a metà di 25 task: mezza lista slitta a destra. \\
18 & OP 10 & Inseriamo l'ID 63 in una lista al limite della capienza: viene allocato nell'ultimo indirizzo di memoria disponibile. \\
19 & OP 10 & Inserimento di un task con P=1 in una lista di venti con P=0: allocazione corretta in coda al blocco. \\
\bottomrule
\end{tabular}
\caption{Ordinamenti e shift.}
\label{tab:test_funz}
\end{table}

\newpage
\subsubsection{Test di Integrazione e Stress}

Qui abbiamo concatenato operazioni in sequenze per vedere se il protocollo START-DONE reggeva.

\begin{table}[H]
\centering
\small
\begin{tabular}{@{}llp{8cm}@{}}
\toprule
\textbf{Test} & \textbf{Operazioni} & \textbf{Condizione analizzata e risultato ottenuto} \\
\midrule
20 & Ins, Inc, Rem, Ins, Inc, Clear & Sequenza di 10 operazioni miste: verifica dello stato della memoria dopo ogni singolo passo. \\
21 & 5 ins + 2 rem + 1 ins + 2 inc & Scenario operativo complesso: l'ordinamento della lista viene mantenuto correttamente fino alla fine. \\
22 & Ins + Rem + Ins + Inc + Rem & Alternanza di inserimenti e rimozioni: verifica puntuale dell'ID restituito a ogni estrazione. \\
23 & 5$\times$(ins+rem) ciclico & Cicli continui di inserimento e rimozione: il modulo gestisce correttamente il passaggio ripetuto tra 1 e 0 elementi. \\
\addlinespace
26 & Qualunque OP & START inviato subito dopo il DONE: il modulo recepisce e avvia immediatamente la nuova operazione. \\
27 & Qualunque OP & Segnale di START mantenuto alto per un tempo prolungato: l'operazione viene eseguita una sola volta, come da protocollo. \\
\bottomrule
\end{tabular}
\caption{Test di integrazione.}
\label{tab:test_integ}
\end{table}

\subsection{Esito della simulazione}

Tutti i 27 test sono passati, per un totale di 78 operazioni.

\begin{table}[H]
\centering
\begin{tabular}{lr}
\toprule
Parametro & Dettaglio \\
\midrule
Piattaforma & GHDL (macOS ARM64) \\
Standard & VHDL 2008 (--std=08 -frelaxed) \\
Scenari & 27 \\
Aree coperte & 8 sottogruppi \\
Operazioni & 78 \\
Tempo simulato & 44.64 $\mu$s \\
Clock & 20 ns \\
Esito & 27/27 \\
\bottomrule
\end{tabular}
\caption{Riepilogo test.}
\label{tab:simulazione}
\end{table}

\newpage   
Il terminale ha chiuso con:

\begin{quote}
\texttt{===== TUTTI I TEST EDGE CASE COMPLETATI CON SUCCESSO =====}
\end{quote}

\subsection{Validazione con la sintesi FPGA}

Abbiamo sintetizzato il VHDL con Xilinx Vivado WebPACK 2018.3 su FPGA \texttt{Artix-7 xc7a200tfbg484-1} a 50\,MHz (20\,ns). Sintesi completata senza errori.

I test della specifica sul sintetizzato hanno dato:
\begin{itemize}
    \item \textbf{Post-sintesi funzionale}: 9/9.
    \item \textbf{Post-sintesi timing}: 9/9.
\end{itemize}

Ecco le risorse stimate:

\begin{table}[H]
\centering
\begin{tabular}{ll}
\toprule
Dato tecnico & Riscontro \\
\midrule
Software di sintesi & Xilinx Vivado WebPACK 2018.3 \\
Piattaforma target & Artix-7 xc7a200tfbg484-1 \\
Timing di sistema & 20\,ns (50\,MHz) \\
Slice LUTs impiegate & 93 su 8000 \\
Slice Registers occupati & 83 su 16000 \\
Bonded IOB & 54 su 150 \\
BUFGCTRL & 1 su 32 \\
WNS (Worst Negative Slack) & inf \\
WHS (Worst Hold Slack) & inf \\
TNS (Total Negative Slack) & 0.000\,ns \\
Endpoint falliti & Nessuno (0 su 255) \\
Stima consumi On-Chip & 3.934\,W \\
\bottomrule
\end{tabular}
\caption{Sintesi su Vivado.}
\label{tab:sintesi}
\end{table}

\newpage
%=======================================================================
\section{In conclusione}
%=======================================================================

Il modulo \texttt{project\_reti\_logiche} fa quello che deve: esegue le quattro operazioni, si sincronizza con l'esterno via START-DONE, e gestisce i casi limite (lista vuota, priorità satura, task con pari priorità).

La FSM a 26 stati in un unico processo sincrono ha tenuto sotto controllo la complessità e reso deterministici gli accessi alla Block RAM.

Per la validazione: 9 test della specifica più 27 test nostri (base, funzionali, integrazione). Tutti passati, dalla pre-sintesi alla post-sintesi timing. La sintesi su Vivado si è chiusa senza errori.

\end{document}